\subsection{Koncentracja miary w wysokich wymiarach}
Ponieważ pracujemy w warunkach klątwy wymiarowości, nasze metody będą opierać się na losowym próbkowaniu punktów testowych z powierzchni sfery. Matematyczny opis tego procesu wymaga zdefiniowania rotacyjnie niezmienniczej, jednostajnej miary powierzchniowej na sferze $\mathbb{S}^{d-1}$, którą oznaczamy jako $\mu_{\mathbb{S}^{d-1}}$.

W wysokich wymiarach miara ta wykazuje niezwykle silną (często wykładniczą) koncentrację, co oznacza, że losowe rzuty układają się w bardzo regularny, dający się opisać sposób. Wynikiem rzutu miary $\mu_{\mathbb{S}^{d-1}}$ na jednostkową kulę $B_{\mathbb{R}^k}$ w podprzestrzeni $k$-wymiarowej jest miara $\mu_k$, wyrażona wzorem $d\mu_k(y) = \frac{\Gamma(d/2)}{\pi^{k/2}\Gamma((d-k)/2)}(1 - \|y\|_{l_2^k}^2)^{\frac{d-2-k}{2}}dy$. Charakterystyka tego rozkładu udowadnia, dlaczego tak "ślepe" losowanie punktów ostatecznie owocuje trafną reprezentacją naszej funkcji.